


	%Formally, 

\section{An EPTAS for Knapsack Cover with Matroid}



In this section we present an EPTAS for KCM. Generally, our approximation scheme defines an additional {\em partition matroid} constraint that bounds the number of items allowed for the solution from each {\em cost class} - items of almost similar costs, defined by rounding over the costs. Then, we find an optimal solution to the induced matroid intersection instance via {\em Lagrangian relaxation} and refined exchange properties of matroids.
 
An important part of our algorithm requires that we take a {\em basis} of a matroid, rather than an independent set. Thus, in the following we show that this does not make the problem any harder. A {\em KCM basis (KCMB)} instance $I = (E,\cI,c,d,D)$ is a KCM instance. A {\em solution} for $I$ (as a KCMB instance) is a basis $S \in \cI, |S| = \max_{S' \in \cI} |S'|$ such that $d(S) \geq D$. The goal is to find a solution $S$ for $I$ with minimum cost $c(S)$. %that is also solution for $I$ being a KCM instance. 

\begin{lemma}
	\label{lem:b}
%	For any $\rho \geq 1$, there is a $\rho$-approximation for \textnormal{KCMB} if and only if there is a $\rho$-approximation for \textnormal{KCM}. Moreover, there is an \textnormal{EPTAS} for \textnormal{KCM} if and only if there is an \textnormal{EPTAS} for \textnormal{KCMB}.
If there is an \textnormal{EPTAS} for \textnormal{KCMB} then there is an \textnormal{EPTAS} for \textnormal{KCM}.
\end{lemma}

\begin{proof}
	Let $I =  (E,\cI,c,d,D)$ be a KCM instance and let $\eps>$ be an error parameter. %Assume that the optimal solution of $I$ contains at least one element, or the problem is trivial. 
	In addition, let $k = \max_{S' \in \cI} |S'|$ be the cardinality of a basis of the matroid $(E,\cI)$ ($k$ is known as the {\em rank} of the matroid). Finally, let $\cA$ be an EPTAS for KCMB. For every $\alpha \in [k]$ define a KCMB instance $I_{\alpha} = (E,\cI_{\leq \alpha},c,d,D)$ such that $$\cI_{\leq \alpha} = \{S \in \cI~|~|S| \leq \alpha\}.$$ 
	Note that a solution for $I_{\alpha}$ contains exactly $\alpha$ elements since it is a basis of the matroid $(E,\cI_{\leq \alpha})$. Moreover, if $S \in \cI$ is a solution for the KCM instance $I$ it is also a solution for the KCMB instance $I_{\alpha}$, where $\alpha = |S|$. 
	On the other hand, for some $\alpha \in [k]$, if $S_{\alpha} \in \cI$ is a solution for $I_{\alpha}$ it is clearly also a solution for $I$ because $\cI_{\leq \alpha} \subseteq \cI$.  
	
	By the above, to obtain an EPTAS for KCM, we compute $\cA$ on each instance $I_{\alpha}$ for all $\alpha \in [k]$ with the error parameter $\eps$. We return the feasible solution for $I$ found through all iterations which has the minimum cost. Let $\alpha \in [k]$ such that there is a solution $S$ of $I$ that satisfies $c(S) = \OPT(I)$; there is such $\alpha$ since $I$ has a solution. Then for this value of $\alpha$, the EPTAS $\cA$ must return a solution $S_{\alpha}$ for $I_{\alpha}$ such that $c(S_{\alpha}) \leq (1+\eps) \cdot \OPT(I_{\alpha}) = (1+\eps) \cdot \OPT(I)$; we conclude that the constructed algorithm returns a $(1+\eps)$-approximate solution. Since $\cA$ is an EPTAS, there is a function $f:\mathbb{R}_{\geq 0} \rightarrow \mathbb{R}_{\geq 0}$ such that the running time of $\cA$ is bounded by $f(\frac{1}{\eps}) \cdot \io$; since we perform $k = \io$ iterations of $\cA$, the running time of the constructed algorithm is bounded by $f(\frac{1}{\eps}) \cdot \io$. Hence, we obtain an EPTAS for KCM. 
\end{proof}







By Lemma~\ref{lem:b} we first find an EPTAS for KCMB. %For the explanations and definitions in this section, consider a KCMB instance $I =  (E,\cI,c,d,D)$ and an error parameter $\eps>0$. 
We first need a constant factor approximation of the optimum cost for KCMB. %$\OPT(I)$; 
For this goal, we convert a PTAS for KCM to a PTAS for KCMB. \begin{lemma}
	\label{lem:B}
	There is an algorithm \textnormal{PTAS-KCMB} that is a \textnormal{PTAS} for \textnormal{KCMB}.
\end{lemma}

\begin{proof}
	let $\cA$ be a PTAS for KCM; such a PTAS exists by \cite{chakaravarthy2013knapsack}. Let $I =  (E,\cI,c,d,D)$ be a KCMB instance and let $\eps>$ be an error parameter. %Assume that $I$ has an optimal solution and that the optimal solution contains at least one element, or the problem is trivial. 
	Let $k = \max_{S' \in \cI} |S'|$ be the rank of the matroid $(E,\cI)$. Define a KCM instance $I' = (E,\cI,c,d',D')$ such that the following holds. \begin{enumerate}
		\item For all $e \in E$ define $d'(e) = d(e)+2 \cdot d(E)$.
		\item $D' = D+k \cdot 2 \cdot d(E)$. 
		%\item $\cI_{\leq \alpha} = \{S \in \cI~|~|S| \leq \alpha\}$. 
	\end{enumerate} Observe that a solution for $I'$ must contain exactly $k$ elements to cover the demand $D'$ and to be independent in the matroid $(E,\cI)$. Thus, $S \subseteq E$ is a solution for (the KCMB instance) $I$ if and only if it is a solution for (the KCM instance) $I'$ due to the new sizes and demand. 
	%
	By the above, we define algorithm PTAS-KCMB by returning the solution that $\cA$ computes for $I'$ and $\eps$. This results in a PTAS for KCMB.
	%
	%construct the KCM instance $I'$, and compute a $\rho$-approximation for KCM on $I$, which yields a $\rho$-approximation for KCMB as well.
\end{proof}

For our algorithm, we partition the element set of the given instance into {\em cost classes}. %Let $n_I = |E|$ be the number of elements in a KCMB instance $I$. Moreover, 
Let $I =  (E,\cI,c,d,D)$ be a KCMB instance and let $S'$ be a solution for $I$ obtained by using the PTAS PTAS-KCMB for KCMB with the error parameter $\eps'  = 2$ %There is a PTAS for KCMB using the PTAS for KCM \cite{chakaravarthy2013knapsack} and 
described in Lemma~\ref{lem:B}; we set $C_{I,\eps} = c(S')$. Note  that %Note that a PTAS for KCM is also a PTAS for KCMB by Lemma~\ref{lem:b}. 
$\OPT(I) \leq C_{I,\eps} \leq 2 \cdot \OPT(I)$.  For some $\eps>0$, define the $0$-th {\em cost class} of $I$ and $\eps$ as the set of elements with relatively small cost $E_0(I,\eps) = \{e \in E~|~ c(e) < \eps \cdot C_{I,\eps}\}$. Second, let $k_{\eps} = \ceil{\frac{1-\eps}{\eps^2}}$. For all $i \in [k_{\eps}]$ define the $i$-th cost class of $I$ and $\eps$ as \begin{equation}
	\label{eq:class}
	E_i(I,\eps) = \big\{e \in E~|~ \left(\eps+(i-1) \cdot \eps^2\right) \cdot C_{I,\eps} \leq c(e) < (\eps+i \cdot \eps^2) \cdot C_{I,\eps} \big\}. 
\end{equation} let $k = k_{\eps}$ and $E_i = E_i(I,\eps) \forall i \in \{0,1,\ldots,k\}$ when clear from the context. Observe that the difference in cost of any two elements from the same cost class is at most $2\eps^{2} \cdot C_{I,\eps}$ which will be used in the analysis. Our algorithm enumerates over all possibilities to the number of elements chosen for the algorithm from each cost class; we describe each option as a vector. Given a KCMB instance $I =  (E,\cI,c,d,D)$, let $n_I = |E|$ be the number of elements in $I$. Then, given $I$ and $\eps>0$, a {\em pattern vectors} of $I$ and $\eps$ is a vector $(n_0, \ldots, n_k) \in \mathbb{N}^{k+1}$ satisfying the following conditions. \begin{enumerate}
	\item $\sum^{k}_{i = 1} n_i \leq \frac{1}{\eps}$.
	\item $\sum^{k}_{i = 0} n_i \leq n_I$.   
\end{enumerate}	\begin{lemma}
	\label{lem:countK}
	For every \textnormal{KCMB} instance $I$ and $\eps>0$ there are $O\left( (k_{\eps}+1)^{\frac{1}{\eps}} \cdot n_I\right)$ pattern vectors of $I$ and $\eps$. 
\end{lemma}

\begin{proof}
	The number of patterns vectors can be bounded by the number of assignments of $0 \leq x \leq \frac{1}{\eps}$ identical balls into $k_{\eps}$ buckets (over all such  value for $x$), and independently choosing a value for $n_0$ in $\{0,1,\ldots,n_I\}$. Thus, the number of options is bounded by $ (k_{\eps}+1)^{\frac{1}{\eps}} \cdot n_I$.  
\end{proof} 

Wee seek pattern vectors that, informally, describe how many elements the optimum takes from each cost class. More concretely, \begin{definition}
	\label{def:good}
 Let $I =  (E,\cI,c,d,D)$ be a KCMB instance and let $\eps>0$. A pattern vector $(n_0, \ldots, n_k)$ of $I$ and $\eps$ is a {\bf good pattern vector} of $I$ and $\eps$ if there is solution $\OPT$ of $I$ that has an optimal cost $c(\OPT) = \OPT(I)$, such that for all $i \in \{0,1,\ldots,k\}$ it holds that $|\OPT \cap E_i(I,\eps)| = n_i$.	
\end{definition} 
%A pattern vector $(n_0, \ldots, n_k)$ of $I$ and $\eps$ is a {\em good pattern vector} of $I$ and $\eps$ if there is an optimal solution $\OPT$ of $I$ such that for all $i \in \{0,1,\ldots,k\}$ it holds that $|\OPT \cap E_i| = n_i$. %, where $\OPT$ is some fixed optimal solution for $I$. 
By Lemma~\ref{lem:countK} we can enumerate over all pattern vectors in time suitable for an EPTAS; one of them is guaranteed to be a good pattern vector.  %Thus, for the simplicity of the presentation, let $P = (n_0, \ldots, n_k)$  be a fixed good pattern vector, which is determined by enumeration in our scheme. 
%For the following, fix a pattern vector $(n_0, \ldots, n_k)$. 
We use a good pattern vector to define the number of elements allowed for the solution from every cost class. This is modeled by the bounds of a partition matroid.  %, used to bound the number of elements selected from each cost class as in the optimum $\OPT$. 
Formally,

\begin{definition}
	\label{def:partitionMatroid}
	  Let $I =  (E,\cI,c,d,D)$ be a KCMB instance, let $\eps>0$, and let $P = (n_0, \ldots, n_k)$ be a pattern vector of $I$ and $\eps$. Define $(E,\cj_{I,\eps}(P))$ as the partition matroid of $I,\eps,P$ such that \begin{equation}
		\label{eq:J}
		\cj_{I,\eps}(P) = \{S \subseteq E~|~\forall i \in \{0,1,\ldots,k\}:~|S \cap E_i| \leq n_i\}. 
	\end{equation} Let $b = \max_{S \in \cI} |S|$ be the rank of $(E,\cI)$ and let $\cG_{I,\eps}(P) = \{S \in \cI \cap \cj_{I,\eps}(P)~|~|S| = b\}$ be all common bases of $\cI$ and $\cj$. 
\end{definition}

%$\sum^{k}_{i = 1} n_i \leq \frac{1}{\eps}$ which. 





%We assume, w.l.o.g, that for all solutions $S,S'$ of $I$ it holds that $|S| = |S'|$; we can indeed make such an assumption if we alter the size function and the demand of $I$ such that a feasible solution for $I$ must be a basis of $(E,\cI)$; 
The above definition of the set of common bases of the given matroid of an instance $I$ with the constructed partition matroid describes a set that we seek a member of this set with minimum {\em lagrangian weight}, that is a function of the cost and size of the given instance. Specifically, our algorithm relies on the following {\em Lagrangian relaxation} of a KCMB instance $I = (E,\cI,c,d,D)$, $\eps>0$, a pattern vector $P$ of $I$ and a parameter $\lambda' \geq 0$: \begin{equation}
	\label{eq:LR}
	LR_{I,\eps,P}(\lambda') = \min_{S \in \cG_{I,\eps}(P)} c(S)+\lambda' \cdot (D-d(S)). 
\end{equation} Clearly, for any $\lambda' \geq 0$ it holds that $LR_{I,\eps,P}(\lambda') \leq \OPT(I)$, and we consider $\lambda$ that gives the tightest such lower bound: \begin{equation}
	\label{eq:Lambda}
	\lambda_{I,\eps,P} = \argmax_{\lambda' \geq 0} LR_{I,\eps,P}(\lambda').
\end{equation} Observe that for every $\lambda' \geq 0 $ we can compute $LR_{I,\eps,P}(\lambda')$ in time $\io$ since it is an instance of matroid intersection basis maximization, known to be polynomially solvable (see, e.g., \cite{schrijver2003combinatorial}). Therefore, $LR_{I,\eps,P}(\lambda')$ (as a function of $\lambda'$) is a minimum of a set of linear functions (e.g., one for
every $S \in \cG_{I,\eps}(P)$); thus, $LR_{I,\eps,P}(\lambda')$ is a concave and piecewise linear function (see, e.g., Fact 1.1 in \cite{gusfield1983parametric}). Therefore, we can compute $\lambda$ in time $\io$ (see \cite{hassin2004efficient} for a more detailed explanation). 

For the following consider again a KCMB instance $I = (E,\cI,c,d,D)$, $\eps>0$, and a pattern vector $P$ of $I$; for simplicity, let $\lambda = 	\lambda_{I,\eps,P} $. Define the cost function of $\lambda$ as $c_{\lambda} : E \rightarrow \mathbb{R}_{\geq 0}$ such that $c_{\lambda}(e) = c(e)-\lambda \cdot d(e)$ for every $e \in E$. In addition, let $$\co_{I,\eps}(P) = \{S \in\cG_{I,\eps}(P)~|~\forall S' \in \cG_{I,\eps}(P) : c_{\lambda}(S) \leq c_{\lambda}(S')\}$$ be all optimal solutions to the Lagrangian relaxation w.r.t. $I,\eps,P$. %the optimal parameter $\lambda$ (and the good vector pattern $P$). 
The next result generalizes the results of \cite{hassin2004efficient,frederickson1989algorithms,ravi1996constrained} to a general matroid and a partition matroid. The proof  follows from Lemma~3 in \cite{hassin2004efficient}, relaying on the properties of the intersection of a general matroid with a partition matroid provided by \cite{brezovec1988matroid}.\footnote{The proof of Lemma~3 in \cite{hassin2004efficient} is stated for graphic matroids along with a partition matroid; however, the proof does not rely on any property of a specific matroid and applies to general matroids.}
% Let $I =  (E,\cI,c,d,D)$ be a KCMB instance, let $\eps>0$, and let $P = (n_0, \ldots, n_k)$ be a pattern vector of $I$ and $\eps$. 
%\begin{lemma}
%	\label{lem:chain}
%	There is a polynomial time algorithm that given $S,S' \in \co$ finds a sequence $S = S_0, S_1, \ldots, S_{\ell} = S' \in \co$, $\ell \in [n]$, such that for all $i \in \{0,1,\ldots,k\}$ and $j \in \{0,1,\ldots,\ell-1\}$ it holds that \begin{equation*}
%		\label{eq:lem} 
%		\left|  (S_j \setminus S_{j+1}) \cap E_i \right|  = \left|  (S_{j+1} \setminus S_{j}) \cap E_i \right| \leq 1.
%	\end{equation*} %$$ \left|  (S_j \setminus S_{j+1}) \cap E_i \right|  = \left|  (S_{j+1} \setminus S_{j}) \cap E_i \right| \leq 1.$$
%	
%	%satisfying the following. \begin{itemize}
%		%			\item For all 
%		%		\end{itemize}
%\end{lemma}


%\begin{lemma}
%	\label{lem:chain}
%	There is an algorithm \textnormal{\textsf{Exchange}} that given a \textnormal{KCMB} instance $I =  (E,\cI,c,d,D)$, $\eps>0$, a pattern vector $P$ of $I$ and $\eps$, and $S,S' \in \co_{I,\eps}(P)$ returns in time $\io$ a set $S'' \in \co_{I,\eps}(P)$ such that $\left|S^{\prime \prime} \backslash S^{\prime}\right|<\left|S \backslash S^{\prime}\right|=x$ and for all $i \in \{0,1,\ldots,k_{\eps}\}$ it holds that \begin{equation*}
%		\label{eq:lem} 
%		\left|  (S'' \setminus S) \cap E_i(I,\eps) \right|  = \left|  (S \setminus S'') \cap E_i(I,\eps) \right| \leq 1.
%	\end{equation*} %$$ \left|  (S_j \setminus S_{j+1}) \cap E_i \right|  = \left|  (S_{j+1} \setminus S_{j}) \cap E_i \right| \leq 1.$$
%	
%	%satisfying the following. \begin{itemize}
%		%			\item For all 
%		%		\end{itemize}
%\end{lemma}
%
%\begin{proof}
%		Let $\lambda = \lambda_{I,\eps,P}$ and let $c_{\lambda}(e) = c(e)-\lambda \cdot d(e)$ for every $e \in E$. %In addition, if $c_{\lambda}(e) = 0$ for all $e \in E$ the define $\eps' = 1$; otherwise, define $\eps' = \frac{\min_{e \in E} c_{\lambda}(e)}{2 \cdot |E|}$. 
%		We consider an altered cost function $c'$ (w.r.t. a small enough $\eps'>0$) defined as follows. 
%	
%	$$
%	c^{\prime}(e)= \begin{cases}\infty & \text { if } e \notin S \cup S^{\prime} \\ c_{\lambda}(e)+\epsilon^{\prime} & \text { if } e \in S \backslash S^{\prime} \\ c_{\lambda}(e)-\epsilon^{\prime} & \text { if } e \in S^{\prime} \backslash S \\ c_{\lambda}(e) & \text { if } e \in S \cap S^{\prime}\end{cases}
%	$$
%
%	As a consequence of Theorem 4 in \cite{brezovec1988matroid}, we can compute in time $\io$ a common basis $S^{\prime \prime}$ of $\cI$ and $\cj_{I,\eps}(P)$ that has a smaller cost than $S$ w.r.t. $c^{\prime}$ and that for all $i \in \{0,1,\ldots, k_{\eps}\}$ it holds that $|(S \setminus S'') \cap E_i| = |(S'' \setminus S) \cap E_i| \leq 1$; therefore, $S^{\prime \prime} \in \mathcal{O}_{I,\eps}(P)$. Observe that $S^{\prime \prime} \subseteq S \cup S^{\prime}$ otherwise the cost of $S''$ w.r.t. $c'$ is $\infty$ in contradiction that $c'(S'')< c'(S)$. Moreover, $S^{\prime \prime} \neq S$ since $c'(S'')< c'(S)$; therefore, by the definition of the cost function $c'$ and since $c'(S'')< c'(S)$, it holds that $\left|S^{\prime \prime} \backslash S^{\prime}\right|<\left|S \backslash S^{\prime}\right|$.
%\end{proof}



\begin{lemma}
	\label{lem:chain}
	There is an algorithm \textnormal{\textsf{Chain}} that given a \textnormal{KCMB} instance $I =  (E,\cI,c,d,D)$, $\eps>0$, a pattern vector $P$ of $I$ and $\eps$, and $S,S' \in \co_{I,\eps}(P)$ returns in time $\io$ a sequence $S = S_0, S_1, \ldots, S_{\ell} = S' \in \co_{I,\eps}(P)$, $\ell \in [n_I]$, such that for all $i \in \{0,1,\ldots,k_{\eps}\}$ and $j \in \{0,1,\ldots,\ell-1\}$ it holds that \begin{equation*}
		\label{eq:lem} 
		\left|  (S_j \setminus S_{j+1}) \cap E_i(I,\eps) \right|  = \left|  (S_{j+1} \setminus S_{j}) \cap E_i(I,\eps) \right| \leq 1.
	\end{equation*} %$$ \left|  (S_j \setminus S_{j+1}) \cap E_i \right|  = \left|  (S_{j+1} \setminus S_{j}) \cap E_i \right| \leq 1.$$
	
	%satisfying the following. \begin{itemize}
		%			\item For all 
		%		\end{itemize}
\end{lemma}




\begin{algorithm}[h]
	\caption{$\textsf{EPTAS}(I = (E, \cI, c,d, D),\eps)$}
	\label{alg:findRep}
	\SetKwInOut{Input}{input}
	\SetKwInOut{Output}{output}
	
%	\Input{A KCMB instance $I$ and error parameter $\eps>0$. }%\ariel{I removed $\alpha$, it was not required. } }
%
%\Output{A representative set $R$ of $I$ and $\eps$.}

%	Compute a $2$-approximation $S^*$ for $I$ using a PTAS for BM  with parameter $\eps' = \frac{1}{2}$.\label{step:REPAPP}
%	
%	Set $\alpha \leftarrow p(S^*)$. \label{step:REPalpha}

%		\If{$|I|<\ell^{(k-1) \cdot \ell}$}{Return $E$.\label{step:ifR}}
%	
%	Initialize $R \leftarrow \emptyset$.


%\If{$\ell^{(k-1) \cdot \ell } \cdot k^2  \cdot \eps^{-2}> |I|$}{Return $E$.\label{step:IF}}
%
%Compute $\alpha \leftarrow \textsf{ApproxBM}(I)$.\label{step:A}

Compute a solution $A$ of $I$. 

\For{\textnormal{pattern vectors $P$ of $I$ and $\eps$} \label{step:REPforr}}{
	
	Compute $S,S' \in \co_{I,\eps}(P)$ such that $d(S) \leq D \leq d(S')$.
	
	Compute $(S_0,\ldots, S_{\ell_P}) \leftarrow \textsf{Chain}(I,\eps,P,S,S')$.

Let $j \in \{0,1,\ldots\ell_P-1\}$ be the first index such that $c(S_j) \leq LR_{I,\eps}(\lambda_{I,\eps,P}) \leq c(S_{j+1})$
	
	\If{$S_{j+1}$ is a solution of $I$ and $c(S_{j+1})< c(A)$}{

Update $A \leftarrow S_{j+1}$. 	

}
}

Return $A$.\label{step:REPret}
\end{algorithm}

%\begin{proof}
%	Algorithm \textsf{Chain} is defined as follows. If $S = S'$ then return $S$. Otherwise, compute algorithm $\textsf{Exchange}$ on $I,\eps,P,S,S'$ to get a set $S''$. Then, return $S$ concatenated by the result from a recursive application of algorithm $\textsf{Exchange}$ on $I,\eps,P,S'',S'$. Observe that $|S'' \setminus S'| < |S \setminus S'|$ by Lemma~\ref{lem:exchange}; thus, the algorithm stops after at most $|E|$ iterations. Since each iteration take $\io$ computation time, the overall running time of algorithm $\textsf{Chain}$ is $\io$. We prove by induction on $|S' \setminus S|$ that algorithm \textsf{Chain} returns the desired sequence. For the base case, if $|S' \setminus S| = 0$, the claim follows by the stopping condition of the algorithm; otherwise, assume that the algorithm the desired sequence for $|S'\setminus S| = x-1$ for some $x \in \mathbb{N}$; then, by Lemma~\ref{lem:exchange}, the assumption of the induction, we immediately have the claim. 
%	
%%	We show that we can find a sequence as described in the statement of the lemma in time $\io$; we construct the sequence inductively on $|S' \setminus S|$. We show that our algorithm returns the desired sequence in length at most $|S' \setminus S|$ and in time $\io$. For $|S' \setminus S| = 1$, our algorithm simply returns the sequence $S,S'$ which satisfies the conditions of the lemma. Assume that the claim holds for $|S' \setminus S| = x-1$ for some $x$; that is, assume that we can compute in time $\io$.
%%	
%%	we now prove it for $|S' \setminus S| = x$. 
%%	
%%	
%%	Let $\lambda = \lambda_{I,\eps,P}$ and let $c_{\lambda}(e) = c(e)-\lambda \cdot d(e)$ for every $e \in E$. We consider an altered cost function $c'$ that considers a small enough $\eps'>0$ defined as follows. 
%%	
%%	$$
%%	c^{\prime}(e)= \begin{cases}\infty & \text { if } e \notin S \cup S^{\prime} \\ c_{\lambda}(e)+\epsilon^{\prime} & \text { if } e \in S \backslash S^{\prime} \\ c_{\lambda}(e)-\epsilon^{\prime} & \text { if } e \in S^{\prime} \backslash S \\ c_{\lambda}(e) & \text { if } e \in S \cap S^{\prime}\end{cases}
%%	$$
%%	%We take $T$ to be the initial solution to the primal algorithm in \cite{brezovec1988matroid}. $T$ is not optimal with respect to $c^{\prime}$ (because $T^{\prime}$ is the unique optimal solution with respect to $c^{\prime}$ ), and therefore the algorithm finds a negative cycle in a graph $S$. 
%%%	We define  a graph $H$ to be used in the algorithm; for the definition of $H$, we use the following auxiliary bipartite directed graph $B=\left(V, V^{\prime}, A\right)$ (see also \cite{brezovec1986two} and \cite{brezovec1988matroid}). For every $e \in S$ we define a vertex in $V$, and for every $e^{\prime} \in E \backslash S$ we have a vertex in $V^{\prime}$. For every $e_i \in S$ and $e_j^{\prime} \in E \backslash S$ we construct an edge $\left(e_i, e_j^{\prime}\right)$ with cost $c\left(e_i, e_j^{\prime}\right)=c^{\prime}\left(e_i\right)-c^{\prime}\left(e_j\right)$ if and only if $(S-e_i)+e'_j \in \cI$ (is independent in the matroid  $(E,\cI)$); moreover, we construct an edge $\left(e_j^{\prime}, e_i\right)$ with cost $c\left(e_j^{\prime}, e_i\right)=0$ if and only if $(S-e_i)+e'_j \in \cj_{I,\eps}(P)$ (is independent in the partition matroid $(E,\cj_{I,\eps}(P))$); i.e., $e_i$ and $e_j^{\prime}$ belong to the same set of the partition in $\cj_{I,\eps}(P)$. Since for $e \notin S \cup S^{\prime}$ it holds that $c^{\prime}(e)=\infty$, a directed cycle in $B$ with negative total cost (a {\em negative cycle}) cannot contain $e$ since it has an infinite cost and there are only a finite number of edges with negative costs which all have finite costs.  Now, define the graph $H = (V,\{v_0,\ldots,v_k\},\mathcal{E})$ as follows. First, $V$ is the set $S$ as defined for the graph $B$. Second, for every $i \in \{0,1,\ldots,k\}$ we have  a vertex $v_i$, such that for all $e \in V$ there is an edge $(v,v_i) \in \mathcal{E}$ if and only if there is $e' \in (E \setminus S) \cap E_{i}(I,\eps)$ such that $(e,e') \in A$ (is an edge in the graph $B$). 
%%%	
%%%	We construct a negative cycle $\cC$ in $H$ with minimum number of edges among all negative cycles in $B$. The cycle $\cC$ can be computed in time $\io$ (see, e.g., \cite{brezovec1988matroid} for more details). % by uniting the set of vertices corresponding to elements from $(E \backslash S) \cap E_i(I,\eps)$ in $V^{\prime}$ to a single vertex $v_i$ for every $i \in [k_{\eps}]$, and removing parallel edges.
%%	%
%%	By Theorem 4 in \cite{brezovec1988matroid}, we can compute in time $\io$ a common basis $S^{\prime \prime}$ of $\cI$ and $\cj_{I,\eps}(P)$ that has a smaller cost than $S$ w.r.t. $c^{\prime}$ and that for all $i \in \{0,1,\ldots, k_{\eps}\}$ it holds that $|(S \setminus S'') \cap E_i| = |(S'' \setminus S) \cap E_i| \leq 1$; therefore, $S^{\prime \prime} \in \mathcal{O}_{I,\eps}(P)$. 
%%	
%%	%Moreover, $S^{\prime \prime}$ can be computed in time $\io$ from $S$ by swapping edges along the negative cycle $\cC$. Suppose that the sequence of vertices along the cycle $\cC$ corresponds to the elements $\left(e_1, e_1^{\prime}, e_2, e_2^{\prime}, \ldots, e_l, e_l^{\prime}, e_1\right)$ such that $e_r \in E_{i_r}(I,\eps)$ for $r \in [\ell]$ and $e'_r \in S$ for $r \in [\ell]$. Since $\cC$ is a negative cycle with minimum number of edges, it clearly holds that $C$ is a simple cycle; therefore, $i_r \neq i_{r^{\prime}}$ for all $r,r^{\prime} \in [\ell], r \neq r^{\prime}$. Thus, $S^{\prime \prime} \in \mathcal{O}_{I,\eps}(P)$. 
%%	Observe that $S^{\prime \prime} \subseteq S \cup S^{\prime}$ otherwise the cost of $S''$ w.r.t. $c'$ is $\infty$ in contradiction that $c'(S'')< c'(S)$. Moreover, and $S^{\prime \prime} \neq S$ since $c'(S'')< c'(S)$; therefore, $\left|S^{\prime \prime} \backslash S^{\prime}\right|<\left|S \backslash S^{\prime}\right|=x$. Thus, by the induction assumption, we can compute the entire sequence from $S$ as desired. Observe that adding one set $S_i$ to the sequence takes $\io$ time by By Theorem 4. Moreover, the length of the sequence is at most $|S \setminus S''| \leq |E|$; thus, algorithm $\textsf{Chain}$ runs in time $\io$. 
%\end{proof}

%\begin{proof}
%	We show that we can find the described sequence in time $\io$; we construct the sequence inductively on $|S' \setminus S|$. For $|S' \setminus S| = 1$, the claim trivially holds and the algorithm just returns the sequence $S,S'$. Assume that the claim holds for $|S' \setminus S| = x-1$ for some $x$; we now prove it for $|S' \setminus S| = x$. Let $\lambda = \lambda_{I,\eps,P}$ and let $c_{\lambda}(e) = c(e)-\lambda \cdot d(e)$ for every $e \in E$. We consider an altered cost function $c'$ that considers a small enough $\eps'>0$ defined as follows. 
%	
%	$$
%	c^{\prime}(e)= \begin{cases}\infty & \text { if } e \notin S \cup S^{\prime} \\ c_{\lambda}(e)+\epsilon^{\prime} & \text { if } e \in S \backslash S^{\prime} \\ c_{\lambda}(e)-\epsilon^{\prime} & \text { if } e \in S^{\prime} \backslash S \\ c_{\lambda}(e) & \text { if } e \in S \cap S^{\prime}\end{cases}
%	$$
%	%We take $T$ to be the initial solution to the primal algorithm in \cite{brezovec1988matroid}. $T$ is not optimal with respect to $c^{\prime}$ (because $T^{\prime}$ is the unique optimal solution with respect to $c^{\prime}$ ), and therefore the algorithm finds a negative cycle in a graph $S$. 
%	We define an auxiliary bipartite directed graph $B=\left(V, V^{\prime}, A\right)$ as follows (see \cite{brezovec1986two} and \cite{brezovec1988matroid}). For every $e \in S$ we define a vertex in $V$, and for every $e^{\prime} \in E \backslash T$ we have a vertex in $V^{\prime}$. For every $e_i \in S$ and $e_j^{\prime} \in E \backslash S$ we construct an edge $\left(e_i, e_j^{\prime}\right)$ with cost $c\left(e_i, e_j^{\prime}\right)=c^{\prime}\left(e_i\right)-c^{\prime}\left(e_j\right)$ if and only if $(S-e_i)+e'_j \in \cI$ (is independent in the matroid  $(E,\cI)$); moreover, we construct an edge $\left(e_j^{\prime}, e_i\right)$ with cost $c\left(e_j^{\prime}, e_i\right)=0$ if and only if $(S-e_i)+e'_j \in \cj_{I,\eps}(P)$ (is independent in the partition matroid $(E,\cj_{I,\eps}(P))$); i.e., $e_i$ and $e_j^{\prime}$ belong to the same set of the partition in $\cj_{I,\eps}(P)$. Since for $e \notin S \cup S^{\prime}$ it holds that $c^{\prime}(e)=\infty$, a directed cycle in $B$ with negative total cost (a {\em negative cycle}) cannot contain $e$ since it has an infinite cost and there are only a finite number of edges with negative costs which all have finite costs.  
%	
%	We construct a negative cycle $\cC$ in $B$. The cycle $\cC$ is obtained from $B$ (see \cite{brezovec1988matroid}) by uniting the set of vertices corresponding to elements from $(E \backslash S) \cap E_i(I,\eps)$ in $V^{\prime}$ to a single vertex $v_i$ for every $i \in [k_{\eps}]$, and removing parallel edges.
%	
%	By Theorem 4 in \cite{brezovec1988matroid}, there is a common base $S^{\prime \prime}$ of $\cI$ and $\cj_{I,\eps}(P)$ that has a smaller cost than $S$ with respect to $c^{\prime}$; therefore, $S^{\prime \prime} \in \mathcal{O}_{I,\eps}(P)$, and it holds that such that $S^{\prime \prime}$ is obtained from $S$ by swapping edges along a negative cycle $\cC$ that has a minimum number of edges (among all the negative cycles in $S$). Suppose that the nodes along the cycle $C$ correspond to edges $\left(e_1, e_1^{\prime}, e_2, e_2^{\prime}, \ldots, e_l, e_l^{\prime}, e_1\right)$ such that $e_k \in E_{i_k}$ for $k=1, \ldots, l$. $C$ is simple, and therefore, $i_k \neq i_{k^{\prime}} \forall k \neq k^{\prime}$ and $T^{\prime \prime} \in \mathcal{O}$.
%	$T^{\prime \prime} \subseteq T \cup T^{\prime}$ and $T^{\prime \prime} \neq T$, and therefore, $\left|T^{\prime \prime} \backslash T^{\prime}\right|<\left|T \backslash T^{\prime}\right|=x$. Therefore, by the induction assumption, the lemma's claim holds.
%\end{proof}


%\begin{proof}
%	We show that we can find the described sequence in time $\io$; we construct the sequence inductively on $|S' \setminus S|$. For $|S' \setminus S| = 1$, the claim trivially holds and the algorithm just returns the sequence $S,S'$. Assume that the claim holds for $|S' \setminus S| = x-1$ for some $x$; we now prove it for $|S' \setminus S| = x$. Let $\lambda = \lambda_{I,\eps,P}$ and let $c_{\lambda}(e) = c(e)-\lambda \cdot d(e)$ for every $e \in E$. We consider an altered cost function $c'$ that considers a small enough $\eps'>0$ defined as follows. 
%	
%	$$
%	c^{\prime}(e)= \begin{cases}\infty & \text { if } e \notin S \cup S^{\prime} \\ c_{\lambda}(e)+\epsilon^{\prime} & \text { if } e \in S \backslash S^{\prime} \\ c_{\lambda}(e)-\epsilon^{\prime} & \text { if } e \in S^{\prime} \backslash S \\ c_{\lambda}(e) & \text { if } e \in S \cap S^{\prime}\end{cases}
%	$$
%	%We take $T$ to be the initial solution to the primal algorithm in \cite{brezovec1988matroid}. $T$ is not optimal with respect to $c^{\prime}$ (because $T^{\prime}$ is the unique optimal solution with respect to $c^{\prime}$ ), and therefore the algorithm finds a negative cycle in a graph $S$. 
%	We define  a graph $H$ to be used in the algorithm; for the definition of $H$, we use the following auxiliary bipartite directed graph $B=\left(V, V^{\prime}, A\right)$ (see also \cite{brezovec1986two} and \cite{brezovec1988matroid}). For every $e \in S$ we define a vertex in $V$, and for every $e^{\prime} \in E \backslash S$ we have a vertex in $V^{\prime}$. For every $e_i \in S$ and $e_j^{\prime} \in E \backslash S$ we construct an edge $\left(e_i, e_j^{\prime}\right)$ with cost $c\left(e_i, e_j^{\prime}\right)=c^{\prime}\left(e_i\right)-c^{\prime}\left(e_j\right)$ if and only if $(S-e_i)+e'_j \in \cI$ (is independent in the matroid  $(E,\cI)$); moreover, we construct an edge $\left(e_j^{\prime}, e_i\right)$ with cost $c\left(e_j^{\prime}, e_i\right)=0$ if and only if $(S-e_i)+e'_j \in \cj_{I,\eps}(P)$ (is independent in the partition matroid $(E,\cj_{I,\eps}(P))$); i.e., $e_i$ and $e_j^{\prime}$ belong to the same set of the partition in $\cj_{I,\eps}(P)$. Since for $e \notin S \cup S^{\prime}$ it holds that $c^{\prime}(e)=\infty$, a directed cycle in $B$ with negative total cost (a {\em negative cycle}) cannot contain $e$ since it has an infinite cost and there are only a finite number of edges with negative costs which all have finite costs.  Now, define the graph $H = (V,\{v_0,\ldots,v_k\},\mathcal{E})$ as follows. First, $V$ is the set $S$ as defined for the graph $B$. Second, for every $i \in \{0,1,\ldots,k\}$ we have  a vertex $v_i$, such that for all $e \in V$ there is an edge $(v,v_i) \in \mathcal{E}$ if and only if there is $e' \in (E \setminus S) \cap E_{i}(I,\eps)$ such that $(e,e') \in A$ (is an edge in the graph $B$). 
%	
%	 We construct a negative cycle $\cC$ in $H$ with minimum number of edges among all negative cycles in $B$. The cycle $\cC$ can be computed in time $\io$ (see, e.g., \cite{brezovec1988matroid} for more details). % by uniting the set of vertices corresponding to elements from $(E \backslash S) \cap E_i(I,\eps)$ in $V^{\prime}$ to a single vertex $v_i$ for every $i \in [k_{\eps}]$, and removing parallel edges.
%	%
%	By Theorem 4 in \cite{brezovec1988matroid}, there is a common basis $S^{\prime \prime}$ of $\cI$ and $\cj_{I,\eps}(P)$ that has a smaller cost than $S$ with respect to $c^{\prime}$; therefore, $S^{\prime \prime} \in \mathcal{O}_{I,\eps}(P)$. Moreover, $S^{\prime \prime}$ can be computed in time $\io$ from $S$ by swapping edges along the negative cycle $\cC$. Suppose that the sequence of vertices along the cycle $\cC$ corresponds to the elements $\left(e_1, e_1^{\prime}, e_2, e_2^{\prime}, \ldots, e_l, e_l^{\prime}, e_1\right)$ such that $e_r \in E_{i_r}(I,\eps)$ for $r \in [\ell]$ and $e'_r \in S$ for $r \in [\ell]$. Since $\cC$ is a negative cycle with minimum number of edges, it clearly holds that $C$ is a simple cycle; therefore, $i_r \neq i_{r^{\prime}}$ for all $r,r^{\prime} \in [\ell], r \neq r^{\prime}$. Thus, $S^{\prime \prime} \in \mathcal{O}_{I,\eps}(P)$. 
%	$T^{\prime \prime} \subseteq T \cup T^{\prime}$ and $T^{\prime \prime} \neq T$, and therefore, $\left|T^{\prime \prime} \backslash T^{\prime}\right|<\left|T \backslash T^{\prime}\right|=x$. Therefore, by the induction assumption, the lemma's claim holds.
%\end{proof}

\comment{
	\begin{proof}
		The proof of 
		
		\comment{the lemma is based on exchanging elements from $S$ until reaching $S'$. The proof uses induction on $x = |S \setminus S'|$. For $x = 0$ the claim trivially follows. Assume that the claim holds for any pair $S,S'$ with $|S \setminus S'| = x-1$. We now prove the claim for $|S \setminus S'| = x$. We define an altered cost function to isolate the specific optimal solution $S'$ of $\co$.  For all $e \in E$ and a sufficiently small $\delta>0$ define \begin{equation}
				\label{eq:c'}
				c'(e) = 
				\begin{cases}
					\infty &  e \notin S \cup S', \\
					c_{\lambda}(e)+\delta &  e \in S \setminus S', \\
					c_{\lambda}(e)-\delta &  e \in S' \setminus S, \\
					c_{\lambda}(e) &  e \in S' \cap S. \\
				\end{cases}
			\end{equation} Note that by \eqref{eq:c'} we have that $S'$ is the unique optimal solution of \eqref{eq:LR} with parameter $\lambda$ with the cost function $c'$ (if $\delta$ is small enough). Thus, $S$ is not optimal w.r.t. $c'$. We define the directed bipartite graph $B= (T,E\setminus T,A)$ (i.e., the vertices of $B$ are $E$, with the bipartition $T,E \setminus T$ and the edge set is $A$) such that the following holds. \begin{enumerate}
				\item $B$
		\end{enumerate}}
\end{proof}}
%	Thus, using the work of Megiddo \cite{megiddo1978combinatorial} we can find $M_1, M_2 \in \cm(G), M_1 \cap L = M_2 \cap L = L^*$ such that the following holds. 
%		
%		\begin{equation}
	%			\label{eq:Megido}
	%			\begin{aligned}
		%				c(M_2) \leq{} & D \leq c(M_1), \\
		%				\\
		%				c_{\lambda} (M_1) ={} & c_{\lambda} (M_2) = LR(\lambda). 
		%			\end{aligned}
	%		\end{equation} We can now state the following inequality, for each matching $M \in \{M_1, M_2\}$: \begin{equation}
	%			\label{eq:opt}
	%			s
	%		\end{equation}



\begin{lemma}
	\label{lem:EPTASbound}
	We can find $S \in \co$ such that $c(S) \leq (1+4 \cdot \eps) \cdot \OPT(I)$ and $d(S) \geq D$. 
\end{lemma}

\begin{proof}
	Since the matroid intersection problem is polynomially solvable, the problems described in \eqref{eq:LR} and \eqref{eq:Lambda} are polynomially solvable. Thus, using the parametric search technique \cite{megiddo1978combinatorial} we can find in polynomial time $S,S' \in \co$ such that $d(S) \leq D \leq d(S')$ and therefore $c(S) \leq LR(\lambda) \leq c(S')$. Then, by Lemma~\ref{lem:chain} we can find a sequence $S = S_0, S_1, \ldots, S_{\ell} = S' \in \co$, $\ell \in [n]$ satisfying the conditions of the lemma. Since $c(S') \leq LR(\lambda) \leq c(S)$, there is $j \in \{0,1,\ldots,\ell-1\}$ such that $c(S_j) \leq LR(\lambda) \leq c(S_{j+1})$. Hence, \begin{equation}
		\label{eq:2}
		c(S_{j+1}) \leq c(S_j)+\sum_{i \in \{0,1,\ldots,k\}} \left(\max_{e \in E_i} c(e) - \min_{e \in E_i} c(e)\right) \leq \OPT(I) + \eps \cdot C+ \frac{1}{\eps} \cdot \eps^2 C \leq (1+4 \cdot \eps) \cdot \OPT(I). 
	\end{equation} The first inequality holds by Lemma~\ref{lem:chain} and \eqref{eq:class}. The second inequality follows since the class $E_0$ may contribute to the sum at most $\eps \cdot C$, and only other $\frac{1}{\eps}$ cost classes can contribute to the sum, where the contribution of such a cost class is bounded by $\eps^2 \cdot C$ due to \eqref{eq:class}. The last inequality holds since $C \leq 2 \cdot \OPT$.  Finally, since $S_{j+1} \in \co$ and $c(S_{j+1} \geq LR(\lambda))$ it follows that $d(S_{j+1}) \geq D$. 
	%			Let $\lambda^{-} = \lambda - \delta$ for a small enough $\delta>0$. It follows that there is $S^- \in \co$ that is optimal also w.r.t. the cost function $c_{\lambda^{-}}$ (rather than $c_{\lambda}$). Since $c_{\lambda}(S^-) = LR(\lambda^-) \leq LR(\lambda)$ it follows that \begin{equation}
		%				\label{eq:-}
		%			c(S^-)+ \lambda^- \cdot (D-d(S^-)) \leq c(S^-) +  \lambda \cdot (D-d(S^-)) = LR(\lambda).
		%			\end{equation} The equality follows since $S^- \in \co$. 
\end{proof}

The above gives an EPTAS for KCMB and therefore also to KCM by Lemma~\ref{lem:b}. 	(TODO: add pseudocode and further details)
%	We can now describe our EPTAS. Given a KCM instance we first reduce it to a family of KCMB instances $I_{\alpha}$, for all $\alpha$ that can be the cardinality of the optimum of $I$ (i.e., $\alpha \in [b]$ where $b$ is the maximal independent set in the given matroid), as in the proof of Lemma~\ref{lem:b}. For each instance $I_{\alpha}$ we compute the two approximation $C$, enumerate over all pattern vectors, for each pattern vector we define a corresponding partition matroid


\comment{
\section{Hardness result for KCM}


(TODO: change the size function to $d$)


%In the following, we define a family of matroids. 
\begin{lemma}
	\label{lem:KCM}
	There is no \textnormal{FPTAS} for \textnormal{KCM} unless \textnormal{P=NP}.%is Strongly \textnormal{NP}-complete
\end{lemma}

\noindent	{\em Proof.}
%Let $X = (E,v,\sigma,H,T)$ be an GSS instance and let $n = |E|$. 
We give a reduction from the $k$-clique problem. Let $G = (V,E)$ be a graph and let $k \in \mathbb{N}_{>0}$; the goal is to decide whether there is a clique of size $k$ in $G$. Let $C(G)$ be all cliques in $G$. Let $|V| = n$ and let $\sigma : V \rightarrow [n]$ be an arbitrary bijeciton. For every $\alpha \in [n^2]$ define the matroid $M_{\alpha} = (V,\cI_{\alpha})$, $\cI_{\alpha} = \cI^1_{\alpha} \cup \cI^2_{\alpha} \cup \cI^3_{\alpha}$, such that \begin{equation}
	\label{eq:1}
	\begin{aligned}
		\cI^1_{\alpha} = {} & \{S \subseteq V~\big|~ |S| < k\} \\
		\cI^2_{\alpha} = {} & \{S \subseteq V~\big|~ |S| = k, \sigma(S) \neq \alpha\} \\
		\cI^3_{\alpha} = {} & \{S \subseteq V~\big|~ |S| = k, \sigma(S) = \alpha, S \in C(G)\}. \\
	\end{aligned}
\end{equation}
\begin{claim}
	\label{claim:1}
	For all $\alpha \in [n^2]$ it holds that $M_{\alpha}$ is a matroid.
\end{claim}
\begin{proof}
	First, note that $\emptyset \in \cI^1_{\alpha}$ since $k \geq 1$ and it follows that $\emptyset \in \cI_{\alpha}$. For the hereditary property, let $A \in \cI_{\alpha}$. For all $B \subset A$ it holds that $|B| < k$; thus, $B \in \cI^1_{\alpha}$ and it follows that $B \in \cI_{\alpha}$. For the exchange property, let $A,B \in \cI_{\alpha}$ such that $|A| > |B|$. We consider the following cases. \begin{itemize}
		\item $|B| < k-1$. Then, for all $e \in A \setminus B$ it holds that $|B+e| < k-1+1 = k$; thus, $B+e \in \cI^1_{\alpha}$ and it follows that $B+e \in \cI_{\alpha}$. Note that there is such $e  \in A \setminus B$ because $|A| > |B|$. 
		
		\item $|B| = k-1$ and $|A| = k$. We consider two sub cases.  \begin{enumerate}
			\item $B \subseteq A$. Then, because $|A| > |B|$ there is $e \in A \setminus B$; thus, $B+e = A$ (because $|B| = k-1$ and $|A| = k$). Hence, since $A \in \cI_k$ it follows that $B+e \in \cI_k$.
			
			\item $B \not\subseteq A$.  Then, since $|B| = k-1$ and $|A| = k$ it follows that $|B \cap A| < |B| = k-1$; thus, $|A \setminus B| = |A|-|A \cap B| > k - (k-1) = 1$. Thus, there are $e,f \in A \setminus B, e\neq f$; since $\sigma$ is a bijection, it holds that $\sigma(e) \neq \sigma(f)$. Therefore, there is $g \in \{e,f\}$ such that $\sigma(B+g) \neq \alpha$. We conclude that $B+g \in \cI^2_{\alpha}$ and it follows that $B+g \in \cI_{\alpha}$. 
		\end{enumerate}
	\end{itemize}
\end{proof}

%Let $\alpha = 2 \cdot n^2$ and $\beta = 2 \cdot v(E)$. We assume w.l.o.g that $n \geq 1$ and that $v(E) > 0$. 
Let $r = 2n^2$. For every ${\alpha} \in  [n^2]$ define the KCM instance $I_{\alpha} = (V,\cI_{\alpha},c,s,D_{\alpha})$ such that the following holds. \begin{equation}
	\label{eq:def}
	\begin{aligned}
		c(e) ={} &  \sigma(e)  ~~\forall e \in V \\
		s(e) ={} & r+\sigma(e)~~\forall e \in V \\
		D_{\alpha} ={} &  k \cdot r+\alpha. %eta \cdot H+T.
	\end{aligned}
\end{equation}

We need to resolve a small technical issue for the single case where $I_{\alpha}$ is not well defined. Since $\sigma$ is a bijection, there is exactly one set $H \subseteq E$ that contains the $k$ elements with maximal values according to $\sigma$. If $\sigma(H) = \alpha$ and $H \notin C(G)$, then there is no solution for $I_{\alpha}$; in this case, we call $I_{\alpha}$ a {\em pseudo} KCM instance. Alternatively, for any $\alpha \in [n^2]$ it holds that $H \in \cI_{\alpha}$ and $s(H) \geq D_{\alpha}$ and thus $I_{\alpha}$ has a solution and is well defined. %If $H \in \cI_{\alpha}$, then we say that $I_{\alpha}$ is a {\em }

%taking the set $H$ of the $k$ vertices with the largest value of $\sigma$ is a solution for $\alpha$: either $\sigma(H) \neq \alpha$ and then $H \in \cI^2_{\alpha}$ and therefore $H \in \cI_{\alpha}$; otherwise, if $\sigma(H) = \alpha$ then we repeat the construction with any bijection $\sigma': V \rightarrow [n]$,
\comment{
	\begin{claim}
		\label{claim:aux1}
		Let $k \in [n]$ and let $S_k$ be an optimal solution for $I_k$. Then the following holds. \begin{enumerate}
			\item $|S_k| \geq k$. 
			\item %If $|S_k| < k$ then $s(S_k)< k \cdot \alpha+H$ and 
			$c(S_k) \geq k \cdot \alpha+H$.
			\item It holds that $\sigma(S_k) \leq H$. 
		\end{enumerate}
	\end{claim}
	
	\begin{proof}
		For the first property of the claim, assume towards a contradiction that $|S_k| < k$. Then,  \begin{equation*}
			\label{eq:sB2}
			\begin{aligned}
				c(S_k) = {} & |S_k| \cdot \alpha+\sigma(S_k) \leq (k-1) \cdot \alpha+\sigma(E) \leq (k-1) \cdot 2 \cdot n^2 +n^2 = k \cdot \alpha-n^2  < k \cdot \alpha \leq k \cdot \alpha +D. 
			\end{aligned}
		\end{equation*} The first inequality holds since $\sigma$ is non negative and assuming that $|S_k| < k$. The second inequality uses that $\sigma(E) \leq n^2$. %Now, if $|S_k| = k$
		For the second property of the claim, assume towards a contradiction that $\sigma(S_k)>D$. Then, \begin{equation*}
			\label{eq:sB2}
			\begin{aligned}
				c(S_k) = {} & \beta \cdot \sigma(S_k)+v(S_k) \geq \beta \cdot (D+1)+v(S_k) \geq \beta \cdot D+\beta = \beta \cdot D+2 \cdot v(E) > \beta \cdot D+T = B.
			\end{aligned}
		\end{equation*} The first inequality follows from our assumption and that $\sigma$ gives integral values for the elements. By the above, we reach a contradiction that $S_k$ is a solution for $I_k$. Therefore, we conclude that $\sigma(S_k) \leq D$ and it follows that $p(S_k) \leq k \cdot \alpha+D$ (because $|S_k| \leq k$).  
	\end{proof}
	
	\begin{claim}
		\label{claim:iff}
		There is a solution for the \textnormal{GSS} instance $X$ if and only if there is $k \in [n]$ such that the optimal solution for $I_k$ has cost exactly $k \cdot \alpha+H$. 
	\end{claim}
	
	\begin{proof}
		For the first direction, let $Y \subseteq E$ be a solution for $X$, i.e., it holds that $v(Y) = T$ and $\sigma(Y) = D$. Let $k = |Y|$.  Observe that $c(Y) = \beta \cdot \sigma(Y) +v(Y) = \beta \cdot D+T = B$; moreover, $|Y| = k ,v(Y) = T$, and $\sigma(Y) = D$ and it follows that $Y \in \cI^3_k \subseteq \cI_k$. Thus, $Y$ is a feasible solution for $I_k$. Moreover, it holds that $p(Y) = \alpha \cdot k+\sigma(Y) = \alpha \cdot k+D$; thus, $Y$ is an optimal solution for $I_k$ by Claim~\ref{claim:aux1}. 
		
		For the second direction, let  $k \in [n]$ and let $S_k$ be an optimal solution for $I_k$ of profit exactly $k \cdot \alpha+D$. By Claim~\ref{claim:aux1} and that $p(S_k) = k \cdot \alpha +D$ it follows that $|S_k| \geq k$; we get that $|S_k| = k$ because $S_k \in \cI_k$. Moreover, it follows that $\sigma(S_k) = D$ to have that $p(S_k) = k \cdot \alpha +D$. Then, because $S_k \in \cI_k, |S_k| = k$, and $\sigma(S_k) = D$ it follows by \eqref{eq:1} that $v(S_k) = T$. Overall, we conclude that $S_k$ is a solution for $X$. %and $\sigma(S_k) = D$.  %, assume that there is no solution for $X$. 
	\end{proof}
	
	%By Claim~\ref{claim:iff}, we can 
	
}

\begin{claim}
	\label{claim:alpha}
	For any $\alpha \in [n^2]$ and a solution $S$ for $I_{\alpha}$ it holds that $c(S) \geq \alpha$. 
\end{claim}

\begin{proof}
	We have
	$$c(S) = \sigma(S) = s(S)-r \cdot |S| \geq D_{\alpha} - r \cdot |S| \geq D_{\alpha} - r \cdot k = \alpha.$$ The equalities follow from \eqref{eq:def}. The inequalities hold since $S$ is a solution for $I_{\alpha}$: the first inequality holds since $s(S) \geq D_{\alpha}$ and the second inequality holds since $S \in \cI_{\alpha}$ and thus $|S| \leq k$. 
\end{proof}



\begin{claim}
	\label{claim:iff}
	There is a a clique of size $k$ in $G$ if and only if there is $\alpha \in [n^2]$ such that there is a solution for $I_{\alpha}$ with cost exactly $\alpha$. 
\end{claim}

\begin{proof}
	For the first direction, let $Q$ be a clique of size $k$ in $G$. Let $\sigma(Q) = \alpha$. Then, \begin{equation*}
		\label{eq:D}
		s(Q) = k \cdot r+\sigma(Q) = k \cdot r+\alpha = D_{\alpha}
	\end{equation*} Moreover, since $|Q| = k$, $\sigma(Q) = \alpha$ and $Q \in C(G)$ it follows that $Q \in \cI^3_{\alpha}$ and therefore $Q \in \cI_{\alpha}$. We conclude that $Q$ is a solution for $I_{\alpha}$ with cost $c(Q) = \sigma(Q) = \alpha$. For the second direction, let $\alpha \in [n^2]$ and let $T$ be a solution for $I_{\alpha}$ with cost exactly $\alpha$. Note that if $|T| < k$ it holds that $$s(T) \leq (k-1) \cdot r+ \alpha \leq (k-1) \cdot r+n^2 < k \cdot r < k \cdot r+\alpha = D_{\alpha}.$$ Since in this case $T$ is not a solution for $I_{\alpha}$ we conclude that $|T| \geq k$; moreover, by \eqref{eq:1} it follows that $|T| = k$ since $T \in \cI_{\alpha}$.  In addition, $c(T) = \sigma(T) = \alpha$; thus, since $T \in \cI_{\alpha}$, $|T| = k$, and $\sigma(T) = \alpha$, we conclude that $T \in \cI^3_{\alpha}$. Therefore, $T \in C(G)$ where $|T| = k$ as required.  
	%	For the first direction, let $Y \subseteq E$ be a solution for $X$, i.e., it holds that $v(Y) = T$ and $\sigma(Y) = D$. Let $k = |Y|$.  Observe that $c(Y) = \beta \cdot \sigma(Y) +v(Y) = \beta \cdot D+T = B$; moreover, $|Y| = k ,v(Y) = T$, and $\sigma(Y) = D$ and it follows that $Y \in \cI^3_k \subseteq \cI_k$. Thus, $Y$ is a feasible solution for $I_k$. Moreover, it holds that $p(Y) = \alpha \cdot k+\sigma(Y) = \alpha \cdot k+D$; thus, $Y$ is an optimal solution for $I_k$ by Claim~\ref{claim:aux1}. 
	%	
	%	For the second direction, let  $k \in [n]$ and let $S_k$ be an optimal solution for $I_k$ of profit exactly $k \cdot \alpha+D$. By Claim~\ref{claim:aux1} and that $p(S_k) = k \cdot \alpha +D$ it follows that $|S_k| \geq k$; we get that $|S_k| = k$ because $S_k \in \cI_k$. Moreover, it follows that $\sigma(S_k) = D$ to have that $p(S_k) = k \cdot \alpha +D$. Then, because $S_k \in \cI_k, |S_k| = k$, and $\sigma(S_k) = D$ it follows by \eqref{eq:1} that $v(S_k) = T$. Overall, we conclude that $S_k$ is a solution for $X$. %and $\sigma(S_k) = D$.  %, assume that there is no solution for $X$. 
\end{proof}

Given an FPTAS for KCM we can decide if there is a clique of size $k$ in $G$. We first consider the only value  $\alpha' \in [n^2]$ such that $I_{\alpha'}$ is possibly a pseudo KCM instance (i.e., may not have a feasible solution). In this case, we just check if the set $H$ of the $k$ items with largest values according to $\sigma$ is a clique in $G$; if $H \in C(G)$, we return that there is a clique of size $k$ in $G$. Otherwise, we iterate over all other values $\alpha \in  [n^2] \setminus \{\alpha'\}$, and compute an optimal solution $S_{\alpha}$ for $I_{\alpha}$ using the FPTAS with parameter $\eps = \frac{1}{2 \cdot n^2}$ (these are all well defined instances). Note that the running time is polynomial since $\frac{1}{\eps}$ is polynomial in the input size. If there is $\alpha \in  [n^2] \setminus \{\alpha'\}$ such that $c(S_{\alpha}) = \alpha$, we return that here is a clique of size $k$ in $G$; otherwise, if for all $\alpha \in  [n^2] \setminus \{\alpha'\}$  it holds that  $c(S_{\alpha}) \neq \alpha$, we return that there is no clique of size $k$ in $G$. Note that $\OPT(I_{\alpha}) \leq n \cdot \max_{e \in E} c(e) = n^2$. Thus, 
by the approximation guarantee of the FPTAS with our parameter $\eps$ it follows that $$c(S_{\alpha}) \leq (1+\eps) \cdot \OPT(I_{\alpha}) \leq \OPT(I_{\alpha}) + \frac{n^2}{2 \cdot n^2} < \OPT(I_{\alpha})+1.$$ Since the costs of all the elements are integers, we conclude that $c(S_{\alpha}) = \OPT(I_k)$. %We use the FPTAS with parameter $\eps = \frac{1}{3 \cdot n^3}$. 
%Since for all $k \in [n]$ it holds that $\OPT(I_k)$

By Claim~\ref{claim:alpha} and Claim~\ref{claim:iff} there is a clique of size $k$ in $G$ if and only if $H \in C(G)$ or there is $\alpha \in [n^2] \setminus \{\alpha'\}$ such that $\OPT(I_{\alpha}) = \alpha$. Therefore, we decide correctly whether there is a clique of size $k$ in $G$ in polynomial time. Since the $k$-clique problem is known to be NP-complete, we conclude that KCM does not have an FPTAS unless P=NP.
$\blacksquare$



\section{A Hardness Result for Constrained Minimum Matroid Basis}

Formally, a {\em  constrained minimum matroid masis (CMB)} instance is a tuple $I =  (E,\cI,c,s,L)$, where $(E,\cI)$ is a matroid along with a membership oracle, $c:E \rightarrow \mathbb{R}_{\geq 0}$  gives the element costs, $s:E \rightarrow \mathbb{R}_{\geq 0}$ gives the elements sizes, and $L \in \mathbb{R}_{\geq 0}$ is a budget. A {\em solution} of $I$ is a basis $B \in \cI$ such that $s(B) \leq L$; in addition, the {\em cost} of the solution is $c(B)$ and let $\OPT(I)$ denote the cost of a solution of $I$ with minimum cost. The goal is to find a solution of $I$ of minimum cost. 


%	\section{Hardness result for KCM}


%	(TODO: change the size function to $d$)

%In the following, we define a family of matroids. 
\begin{lemma}
	\label{lem:KCM2}
	There is no \textnormal{FPTAS} for \textnormal{CMB} unless \textnormal{P=NP}.%is Strongly \textnormal{NP}-complete
\end{lemma}

\noindent	{\em Proof.}
%Let $X = (E,v,\sigma,H,T)$ be an GSS instance and let $n = |E|$. 

We give a reduction from the $k$-clique problem. Let $G = (V,E)$ be a graph and let $k \in \mathbb{N}_{>0}$; the goal is to decide whether there is a clique of size $k$ in $G$. Let $C(G)$ be all cliques in $G$. Let $|V| = n$ and let $\sigma : V \rightarrow [n]$ be an arbitrary bijection. For every $\alpha \in [n^2]$ we consider the matroid $M_{\alpha}$ as defined in the previous section. 

%= (V,\cI_{\alpha})$, $\cI_{\alpha} = \cI^1_{\alpha} \cup \cI^2_{\alpha} \cup \cI^3_{\alpha}$, such that \begin{equation}
	%	\label{eq:1}
	%	\begin{aligned}
		%		\cI^1_{\alpha} = {} & \{S \subseteq V~\big|~ |S| < k\} \\
		%		\cI^2_{\alpha} = {} & \{S \subseteq V~\big|~ |S| = k, \sigma(S) \neq \alpha\} \\
		%		\cI^3_{\alpha} = {} & \{S \subseteq V~\big|~ |S| = k, \sigma(S) = \alpha, S \in C(G)\}. \\
		%	\end{aligned}
	%\end{equation}
	%\begin{claim}
	%	\label{claim:1}
	%	For all $\alpha \in [n^2]$ it holds that $M_{\alpha}$ is a matroid.
	%\end{claim}
	%\begin{proof}
	%	First, note that $\emptyset \in \cI^1_{\alpha}$ since $k \geq 1$ and it follows that $\emptyset \in \cI_{\alpha}$. For the hereditary property, let $A \in \cI_{\alpha}$. For all $B \subset A$ it holds that $|B| < k$; thus, $B \in \cI^1_{\alpha}$ and it follows that $B \in \cI_{\alpha}$. For the exchange property, let $A,B \in \cI_{\alpha}$ such that $|A| > |B|$. We consider the following cases. \begin{itemize}
		%		\item $|B| < k-1$. Then, for all $e \in A \setminus B$ it holds that $|B+e| < k-1+1 = k$; thus, $B+e \in \cI^1_{\alpha}$ and it follows that $B+e \in \cI_{\alpha}$. Note that there is such $e  \in A \setminus B$ because $|A| > |B|$. 
		%		
		%		\item $|B| = k-1$ and $|A| = k$. We consider two sub cases.  \begin{enumerate}
			%			\item $B \subseteq A$. Then, because $|A| > |B|$ there is $e \in A \setminus B$; thus, $B+e = A$ (because $|B| = k-1$ and $|A| = k$). Hence, since $A \in \cI_k$ it follows that $B+e \in \cI_k$.
			%			
			%			\item $B \not\subseteq A$.  Then, since $|B| = k-1$ and $|A| = k$ it follows that $|B \cap A| < |B| = k-1$; thus, $|A \setminus B| = |A|-|A \cap B| > k - (k-1) = 1$. Thus, there are $e,f \in A \setminus B, e\neq f$; since $\sigma$ is a bijection, it holds that $\sigma(e) \neq \sigma(f)$. Therefore, there is $g \in \{e,f\}$ such that $\sigma(B+g) \neq \alpha$. We conclude that $B+g \in \cI^2_{\alpha}$ and it follows that $B+g \in \cI_{\alpha}$. 
			%		\end{enumerate}
		%	\end{itemize}
	%\end{proof}
	
	%Let $\alpha = 2 \cdot n^2$ and $\beta = 2 \cdot v(E)$. We assume w.l.o.g that $n \geq 1$ and that $v(E) > 0$. 
	Let $r = 2n^2$ and let $z = \sum^{k}_{i = 2} i $. %and note that $z \in [n^2]$. 
	For every ${\alpha} \in  [n^2] \setminus [z]$ define the CMB instance $I_{\alpha} = (V,\cI_{\alpha},c,s,L_{\alpha})$ such that the following holds. \begin{equation}
		\label{eq:def2}
		\begin{aligned}
			c(e) ={} &  r-\sigma(e)  ~~\forall e \in V \\
			s(e) ={} & \sigma(e)~~\forall e \in V \\
			L_{\alpha} ={} &  \alpha. 
		\end{aligned}
	\end{equation}
	
	We need to resolve a small technical issue for the single case where $I_{\alpha}$ is not well defined. Since $\sigma$ is a bijection, there is exactly one set $H \subseteq E$ that contains the $k$ elements with minimal values according to $\sigma$. If $\sigma(H) = \alpha$ and $H \notin C(G)$, then there is no solution for $I_{\alpha}$ (we either violate the matroid constraint or the knapsack constraint). In this case, we call $I_{\alpha}$ a {\em pseudo} CMB instance. 
	Alternatively, for any $\alpha \in [n^2] \setminus [z]$ it holds that $H \in \cI_{\alpha}$ and $s(H) \leq L_{\alpha}$ and thus $I_{\alpha}$ has at least one solution and is well defined. %If $H \in \cI_{\alpha}$, then we say that $I_{\alpha}$ is a {\em }
	
	
	
	
	
	\begin{claim}
		\label{claim:alpha2}
		For any $\alpha \in [n^2] \setminus [z]$ and a solution $S$ for $I_{\alpha}$ it holds that $c(S) \geq r \cdot k-\alpha$. 
	\end{claim}
	
	\begin{proof}
		We have
		$$c(S) = r \cdot |S|-\sigma(S)  = r \cdot k-\sigma(S) = r \cdot k-s(S) \geq r \cdot k -L_{\alpha} = r \cdot k -\alpha.$$ The equalities follow from \eqref{eq:def2}; also note that $|S| = k$ because $S$ is a basis of $M_{\alpha}$. The inequality holds since $S$ is a solution for $I_{\alpha}$ and it holds that $s(S) \leq L_{\alpha}$. 
	\end{proof}
	
	
	
	\begin{claim}
		\label{claim:iff2}
		There is a a clique of size $k$ in $G$ if and only if there is $\alpha \in [n^2] \setminus [z]$ such that there is a solution for $I_{\alpha}$ with cost exactly $r \cdot k-\alpha$. 
	\end{claim}
	
	\begin{proof}
		For the first direction, let $Q$ be a clique of size $k$ in $G$. Let $\alpha  = \sigma(Q)$ and observe that $\alpha \in [n^2] \setminus [z]$. Then, \begin{equation*}
			\label{eq:D2}
			s(Q) =\sigma(Q) = \alpha = L_{\alpha}. 
		\end{equation*} Moreover, since $|Q| = k$, $\sigma(Q) = \alpha$ and $Q \in C(G)$ it follows that $Q \in \cI^3_{\alpha}$ and therefore $Q \in \cI_{\alpha}$. We conclude that $Q$ is a solution for $I_{\alpha}$ with cost $$c(Q) = r \cdot |Q| - \sigma(Q) = r \cdot k - \sigma(Q) = r \cdot k - \alpha.$$ 
		
		For the second direction, let $\alpha \in [n^2] \setminus [z]$ and let $T$ be a solution for $I_{\alpha}$ with cost exactly $r \cdot k - \alpha$. Since $T$ is a solution for $I_{\alpha}$ it is a basis of $M_{\alpha}$ and it follows that $|T| = k$. 
		%
		%Note that if $|T| < k$ it holds that $$s(T) \leq (k-1) \cdot r+ \alpha \leq (k-1) \cdot r+n^2 < k \cdot r < k \cdot r+\alpha = D_{\alpha}.$$ Since in this case $T$ is not a solution for $I_{\alpha}$ we conclude that $|T| \geq k$; moreover, by \eqref{eq:1} it follows that $|T| = k$ since $T \in \cI_{\alpha}$.  
		%
		In addition, $$\sigma(T) = r \cdot |T|-c(T) = r \cdot k- c(T) = r \cdot k- (r \cdot k - \alpha) = \alpha.$$ Thus, since $T \in \cI_{\alpha}$, $|T| = k$, and $\sigma(T) = \alpha$, we conclude that $T \in \cI^3_{\alpha}$. Therefore, $T \in C(G)$ where $|T| = k$ as required.  
		%	For the first direction, let $Y \subseteq E$ be a solution for $X$, i.e., it holds that $v(Y) = T$ and $\sigma(Y) = D$. Let $k = |Y|$.  Observe that $c(Y) = \beta \cdot \sigma(Y) +v(Y) = \beta \cdot D+T = B$; moreover, $|Y| = k ,v(Y) = T$, and $\sigma(Y) = D$ and it follows that $Y \in \cI^3_k \subseteq \cI_k$. Thus, $Y$ is a feasible solution for $I_k$. Moreover, it holds that $p(Y) = \alpha \cdot k+\sigma(Y) = \alpha \cdot k+D$; thus, $Y$ is an optimal solution for $I_k$ by Claim~\ref{claim:aux1}. 
		%	
		%	For the second direction, let  $k \in [n]$ and let $S_k$ be an optimal solution for $I_k$ of profit exactly $k \cdot \alpha+D$. By Claim~\ref{claim:aux1} and that $p(S_k) = k \cdot \alpha +D$ it follows that $|S_k| \geq k$; we get that $|S_k| = k$ because $S_k \in \cI_k$. Moreover, it follows that $\sigma(S_k) = D$ to have that $p(S_k) = k \cdot \alpha +D$. Then, because $S_k \in \cI_k, |S_k| = k$, and $\sigma(S_k) = D$ it follows by \eqref{eq:1} that $v(S_k) = T$. Overall, we conclude that $S_k$ is a solution for $X$. %and $\sigma(S_k) = D$.  %, assume that there is no solution for $X$. 
	\end{proof}
	
	Given an FPTAS for CMB we can decide if there is a clique of size $k$ in $G$. We first consider the only value $\alpha' = z+1$ such that $I_{\alpha'}$ is possibly a pseudo CMB instance (i.e., may not have a feasible solution). In this case, we just check if the set $H$ of the $k$ items with the smallest values according to $\sigma$ is a clique in $G$; if $H \in C(G)$, we return that there is a clique of size $k$ in $G$. Otherwise, we iterate over all other values $\alpha \in  [n^2] \setminus ([z] \cup \{\alpha'\})$, and compute an optimal solution $S_{\alpha}$ for $I_{\alpha}$ using the FPTAS with parameter $\eps = \frac{1}{2 \cdot n^2}$ (these are all well defined instances). Note that the running time is polynomial since $\frac{1}{\eps}$ is polynomial in the input size. If there is $\alpha \in  [n^2] \setminus ([z] \cup \{\alpha'\})$ such that $c(S_{\alpha}) = r \cdot k-\alpha$, we return that here is a clique of size $k$ in $G$; otherwise, if for all $\alpha \in  [n^2]\setminus ([z] \cup \{\alpha'\})$  it holds that  $c(S_{\alpha}) \neq r \cdot k-\alpha$, we return that there is no clique of size $k$ in $G$. Note that $\OPT(I_{\alpha}) \leq n \cdot \max_{e \in E} c(e) = n^2$. Thus, 
	by the approximation guarantee of the FPTAS with our parameter $\eps$ it follows that $$c(S_{\alpha}) \leq (1+\eps) \cdot \OPT(I_{\alpha}) \leq \OPT(I_{\alpha}) + \frac{n^2}{2 \cdot n^2} < \OPT(I_{\alpha})+1.$$ Since the costs of all the elements are integers, we conclude that $c(S_{\alpha}) = \OPT(I_k)$. %We use the FPTAS with parameter $\eps = \frac{1}{3 \cdot n^3}$. 
	%Since for all $k \in [n]$ it holds that $\OPT(I_k)$
	
	By Claim~\ref{claim:alpha2} and Claim~\ref{claim:iff2} there is a clique of size $k$ in $G$ if and only if $H \in C(G)$ or there is $\alpha \in [n^2] \setminus ([z] \cup \{\alpha'\})$ such that $\OPT(I_{\alpha}) = r \cdot k-\alpha$. Therefore, we decide correctly whether there is a clique of size $k$ in $G$ in polynomial time. Since the $k$-clique problem is known to be NP-complete, we conclude that CMB does not have an FPTAS unless P=NP.
	$\blacksquare$
	
	
	
}






%\begin{tikzpicture}
%	\fill[gray, opacity=0.5] (0,0) rectangle (4,2);
%\end{tikzpicture}

% \begin{figure} %\hspace*{6cm}
%	%	\label{fig:2}
%	\begin{tikzpicture}[ultra thick]
%		\node[draw, rectangle, minimum width=5cm, minimum height=4cm] (R) at (0,0) {};
%		
%		\node[draw=none] at (4.8, 0) {Representative Set $\bf R$};
%		
%		\node[draw=none] at (-4, -0.5) {Exchange Set}; 
%		
%		\node[draw=none] at (0, 1.5) {$\bf {R \cap  \textcolor{teal}{\mathcal{K}_1(\alpha)}}$};
%		\node[draw=none] at (0, 0.5) {$\bf R \cap  \textcolor{red}{\mathcal{K}_2(\alpha)}$};
%		\node[draw=none] at (0, -0.5) {$\bf R \cap \textcolor{blue}{\mathcal{K}_3(\alpha)}$};
%		\node[draw=none] at (0, -1.5) {$\bf R \cap  \textcolor{brown}{\mathcal{K}_4(\alpha)}$};
%		
%		\draw ([yshift=-1cm]R.north west) -- ([yshift=-1cm]R.north east) 
%		;
%		\draw ([yshift=-2cm]R.north west) -- ([yshift=-2cm]R.north east) 
%		;
%		\draw ([yshift=+1cm]R.south west) -- ([yshift=+1cm]R.south east) 
%		;
%		\draw [decorate,
%		decoration = {calligraphic brace}] (-2.7,-1) --  (-2.7,0);
%		
%		\draw [decorate,
%		decoration = {calligraphic brace}] (2.8,2) --  (2.8,-2);
%		% \draw [decorate, decoration={brace, amplitude=10pt, mirror}, xshift=-6pt, yshift=0pt];
%		%([xshift=-6pt]R.south east) -- ([xshift=-6pt]R.north east);
%	\end{tikzpicture}
%	\caption{\label{fig:2} An illustration of our construction of a representative set $R$, using a union of exchange sets, one for each profit class $\mathcal{K}_1(\alpha)$, $\mathcal{K}_2(\alpha)$, $\mathcal{K}_3(\alpha)$, $\mathcal{K}_4(\alpha)$. }
%\end{figure}





%	\begin{figure}
%	\centering
%	\vspace{-4cm} 
%	\begin{tikzpicture}[scale=1.4, every node/.style={draw, circle, inner sep=1pt}]
%		\node (11) at (5.5,4) {$\bf \textcolor{teal}{1}$};
%		\node (22) at (7,4) {$\bf \textcolor{black}{2}$};
%		\node (33) at (5.5,3) {$\bf \textcolor{red}{3}$};
%		\node (44) at (7,3) {$\bf \textcolor{blue}{4}$};
%		\draw (11) -- (22);
%		\draw (11) -- (33);
%		\draw (11) -- (44);
%		\draw (22) -- (44);
%		\draw (33) -- (44);
%		
%		\node[draw=none] at (9, 3.5) {$~\Pi = \{\textnormal{independent sets of } G\}$};
%		
%		\node[draw=none] at (4.5, 3.5) {A graph $G$};
%		% Formula
%		\node[draw=none] at (6.75, 2.5) {$\mathcal{J}_{4,2} = \{\emptyset, \{$\bf \textcolor{teal}{1}$\},\{$\bf \textcolor{black}{2}$\},\{$\bf \textcolor{red}{3}$\},\{$\bf \textcolor{blue}{4}$\}\}, ~~\mathcal{K}_{4,2,5} = \{\{$\bf \textcolor{black}{2}$,$\bf \textcolor{blue}{4}$\},\{$\bf \textcolor{black}{2}$,$\bf \textcolor{teal}{1}$\},\{$\bf \textcolor{blue}{4}$,$\bf \textcolor{red}{3}$\},\{$\bf \textcolor{red}{3}$,$\bf \textcolor{teal}{1}$\}\}, ~~\mathcal{L}_{4,2,5}(\Pi) = \{\{$\bf \textcolor{black}{2}$,$\bf \textcolor{red}{3}$\}\}$};
%	\end{tikzpicture}
%	\vspace{-7.5cm} 
%	\caption{\label{fig:Pi} The independent sets of the $\Pi$-matroid $M_{n,k,\alpha}(\Pi)$, with parameters $n = 4$, $k = 2$, $\alpha = 5$. The secret set $\Pi$ contains all independent sets in the graph $G$, where $\{2,3\}$ is the only independent set in $G$ with $k$ elements.}
%	
%\end{figure}



